%=====================================================================
\section{The theta-lift family: tensor structure}\label{sec:family}
%=====================================================================
For $s\ge5$ let $A_s$, $Q_s$, $a$, $c_0$ and $\gamma$ be as in \S\ref{sec:notation}, and let $u_{\alpha,\beta}$ be the theta lifts of \S\ref{sec:lift} built from
$A_s$ (for $s=7$ these are the $140$ functions of Theorem~\ref{thm:F}).

\begin{remark}[The $S$ versus $S^{-1}$ label convention]\label{rem:Sconv}
Let $S_{\rm lab}$ be the label matrix of Theorem~\ref{thm:F}, i.e.\ $u(-1/w)=S_{\rm lab}u(w)$. For $s=7,8,11$ one checks exactly: $S_{\rm lab}^2=P$,
$(S_{\rm lab}T)^3=I$ on all labels, $(S_{\rm lab}T)^3\neq S_{\rm lab}^2$, and $(S_{\rm lab}^{-1}T)^3=S_{\rm lab}^{-2}$. So the representation of $\mathrm{SL}_2(\Z)$ in
the standard presentation $S^4=1$, $(ST)^3=S^2$ is $\rho(S):=S_{\rm lab}^{-1}$, $\rho(T):=T$. On $W_P$ one has $S_{\rm lab}=S_{\rm lab}^{-1}$. (Category A, exact.)
\end{remark}

\begin{remark}[Label operators]\label{rem:conv}
Label maps (negation, $\gamma$, the local reflections $Z_p$) are realised as monomial matrices $Y$ with $Y[i,j(i)]=\e(\mathrm{ph}_i)$, the convention of the
permutation matrix of the family. With this convention $\mathrm{neg}=\rho(-I)=S^2$, and $\mathrm{neg}$, $\gamma$, $Z_p$ ($p\mid s$) and $Z_q$ ($q\mid s-2$) commute exactly
with $S$ and $T$ (exact exponent arithmetic modulo $2D^2$ at $s=5,7,9,11,13,15,25$). (Category A.)
\end{remark}

\subsection{Exact tensor factorisation at $s=7$}
Let $g=(14,6)/70$, a generator of $L'/\Z^2\cong\Z/35$, and $a_2=(\tfrac12,0)$, so that $a=a_2+a_{\rm odd}$ with $a_{\rm odd}=(\tfrac15,-\tfrac15)\in L'$. In the
adapted basis $u'_{x,\epsilon,\delta}:=u_{xg+\epsilon a_2,\,-\delta a_2}$ ($x\in\Z/35\cong\Z/5\times\Z/7$ by CRT, $\epsilon,\delta\in\{0,1\}$) put
$F_p[x,y]=\e(4xy/p)$, $T_p=\mathrm{diag}\,\e(2x^2/p)$, and let $S_2$, $T_2$ be the $4\times4$ monomial matrices on the theta characteristics $(\epsilon,\delta)$:
$S_2:00\mapsto00$, $01\leftrightarrow10$, $11\mapsto-11$; $T_2:00\leftrightarrow01$, $10\mapsto i\cdot10$, $11\mapsto i\cdot11$.
\begin{theorem}[Tensor factorisation, $s=7$]\label{thm:tensor7}
(i) In the adapted basis $u'_{x,\epsilon,\delta}$ built on the generator $g=(14,6)/70=(7,3)/35$, exactly $S_{\rm lab}=(F_5/\sqrt5)\otimes(F_7/\sqrt7)\otimes S_2$ and $T=T_5\otimes T_7\otimes T_2$.
(ii) $\rho=\rho_5\otimes\rho_7\otimes\rho_2$, where $\rho_p$ is the Weil representation of $(\Z/p,q_p)$ with $q_5=2x^2/5$ and $q_7=x^2/7$, and
$\rho_2=(\e(3/4)S_2^{-1},T_2)$ acting on the $4$-dimensional space of theta characteristics, is a faithful representation of $\mathrm{SL}_2(\Z/4)$ (image of order $48$). Hence $\rho$ factors through
$\mathrm{SL}_2(\Z/140)=\mathrm{SL}_2(\Z/4)\times\mathrm{SL}_2(\mathbb F_5)\times\mathrm{SL}_2(\mathbb F_7)$.
(iii) $\rho_2=A_3\oplus B_1$, where $A_3$ is the span of the three even characteristics $00,01,10$ (irreducible, $T$-spectrum $\{1,-1,i\}$) and $B_1$ is the odd
characteristic $11$ ($S,T\mapsto\e(1/4)$).
(iv) $P=\rho(-I)=Z_5\otimes Z_7\otimes I_4$, where $Z_p$ is the reflection $x\mapsto-x$. Hence
$W_P=(W_5^+\otimes W_7^+\oplus W_5^-\otimes W_7^-)\otimes\C^4$; the mixed local parities are absent. The restrictions are $16W_5^++12W_5^-$,
$12W_7^++8W_7^-$ and $18(A_3+B_1)$.
(v) $\mathrm{tr}\,\rho_5(S)|_{W_5^\pm}=(-1,0)$, $\mathrm{tr}\,\rho_7(S)|_{W_7^\pm}=(0,-1)$, $\mathrm{tr}\,\rho_2(S)|_{A_3}=-i$, $\mathrm{tr}\,\rho_2(S)|_{B_1}=i$, so $\mathrm{tr}\,S=0$ on each of the
four sectors below and $n_\pm=36/36$ on $W_P$.
\end{theorem}
\begin{proof}[Proof outline]
Every entry of $S_{\rm lab}$ is one term $\e(r)/\sqrt{35}$ with $r\in\Q$, and $T$ is monomial, so (i) is a finite identity over $\Q(\zeta_{140})$, checked
exactly. For (ii), $(\overline{F_p}T_p)^3=c\,Z_p$ exactly in $\Z[\zeta_p]$ with $c^2=125$ ($p=5$) and $c^2=-343$ ($p=7$); with Gauss's sign of the quadratic
Gauss sum this gives the scalars $\gamma_5=-1$, $\gamma_7=\e(3/4)$, so $\rho_p$ is exactly the Weil representation of $q_p$; by the congruence theorem for Weil
representations of finite quadratic modules of odd level \cite{Weil}, $\rho_p$ factors through $\mathrm{SL}_2(\Z/p)$. $\rho_2$ satisfies $S^4=I$,
$(ST)^3=S^2=-I$, $T^4=I$ exactly, and the pair group $(\rho_2(g),g\bmod4)$ has exactly $48$ elements. The scalars cancel in the product, which equals
$(S_{\rm lab}^{-1},T)$. (iv) and (v) follow from the local Gauss sums.
\end{proof}
\catA{Exact (rational and cyclotomic arithmetic), conditional on the cited congruence theorem for the odd local factors \cite{Weil} (hypotheses: exact
match with the Weil formulas, odd level). The local forms are determined as follows: a $140$th-root exponent $a$ decomposes as $a=35x_4+28x_5+20x_7$, i.e.\
$x_4\equiv3a\bmod4$, $x_5\equiv2a\bmod5$, $x_7\equiv6a\bmod7$, which gives $q_5=2x^2/5$ and $q_7=x^2/7$. (Since $2$ is a non-square mod $5$ and $3$ is a
non-square mod $7$, these are not isometric to $x^2/5$ and $3x^2/7$.)}

\begin{proposition}[$V$ is the matched-sector space; local reflections]\label{prop:matched}
Let $X^\pm=W_5^\pm\otimes W_7^\pm$. At $s=7$:
(i) negation acts on $u'$ as $Z_5\otimes Z_7\otimes I_4$, and the unit $\gamma$ ($\varepsilon=6+\sqrt{35}$) acts as $I_5\otimes Z_7\otimes\mathrm{diag}(1,1,1,-1)$;
(ii) $V=\mathrm{Fix}(N)\cap\mathrm{Fix}(\gamma)=A_3\otimes X^+\oplus B_1\otimes X^-$, of dimension $36+6=42$, and the lift kills
$V^\perp=B_1\otimes X^+\,(12)\oplus A_3\otimes X^-\,(18)$;
(iii) $V_{\rm Eis}=A_3\otimes X^+$ and the cuspidal part of $V$ is $V_{\rm cusp}=B_1\otimes X^-$. The $24$ cusp components of \S\ref{sec:split} are exactly the nonzero
labels with $(\epsilon,\delta)=(1,1)$, and the $11$ identically vanishing labels are exactly the $6$ inconsistent orbits ($x_5x_7=0$, $\epsilon=\delta=1$);
(iv) each of the $42$ classes is a free $\langle N,\gamma\rangle$-orbit (sizes $1,2,4$: $3,15,24$ classes), with proportionality constants $\pm1$ in the $u'$ basis;
(v) the local reflections $Z_5$, $Z_7$ and $Z_2:=\mathrm{diag}(1,1,1,-1)$ commute with $\rho$. Their signs $(Z_5,Z_7,Z_2)$ are $(+,+,+)$ on $V_{\rm Eis}$, $(-,-,-)$ on $V_{\rm cusp}$,
$(+,+,-)$ on $B_1\otimes X^+$ and $(-,-,+)$ on $A_3\otimes X^-$. Thus $Z_5Z_7=+1$ on $W_P$ (this is $P$) and $Z_7Z_2=+1$ exactly on $V$ (this is $\gamma$).
$S$ commutes with all three, and $n_\pm=d/2$ on every joint eigenspace.
\end{proposition}
\begin{proof}[Proof outline]
(i) The identities $u_{-\alpha,-\beta}=u_{\alpha,\beta}$ and $u_{\gamma\alpha,\gamma\beta}=u_{\alpha,\beta}$ hold as identities of functions, and their actions on $u'$ are computed exactly.
(ii) Hence $V\subseteq\mathrm{Fix}(N)\cap\mathrm{Fix}(\gamma)$, which has dimension $42$. The exact Fourier matrix (sums of $4900$th roots of unity, $|Q|\le25$) has rank $42$
modulo a prime above $\ell=29401\equiv1\bmod4900$; rank cannot increase under reduction, so $\dim V\ge42$. (iii) The constant-term vector is
$\delta_{x=0}\otimes(e_{00}+e_{01})$, a pure tensor with $2$-adic part in $A_3$; its cyclic span is $A_3\otimes X^+$ (the argument is that of Theorem~\ref{thm:dims}(iii)).
The cuspidal functions of the span are those orthogonal to $V_{\rm Eis}$ (the $\lambda=\tfrac14$ constant-term criterion); by unitarity this is $B_1\otimes X^-$.
\end{proof}
\catA{Exact: (i), (ii), (iv), (v) and the label statements of (iii). The mechanism that kills $V^\perp$ is the $\gamma$ phase flip: the unit acts by $-1$
at $7$ and by $-1$ on the odd characteristic, the analogue of the exact vanishing of the $11$ labels (\S\ref{sec:dihedral}). Part (iii) uses the cited
congruence theorem and the standard $\lambda=\tfrac14$ constant-term criterion.}
\catB{Calling $Z_5,Z_7,Z_2$ the Atkin--Lehner operators $W_5,W_7,W_4$ on the lifted forms (not checked on the forms). The Hecke-character reading: a
character contributes only if $\psi(\varepsilon)=\psi(-1)=1$; genus characters are even at the ramified primes and land in $A_3\otimes X^+$; characters odd at
$5$ and $7$ are not genus characters, hence not $\sigma$-invariant, and give Maass cusp forms in $B_1\otimes X^-$. The identification of
$V_{\rm cusp}$ with specific dihedral Hecke eigenforms rests on the floating fits of \S\ref{sec:split}.}

\subsection{General odd $s$}
\begin{theorem}[Tensor factorisation for odd $s$]\label{thm:tensorodd}
Let $s\ge5$ be odd and $n=n_s$. Then $L'/\Z^2\cong\Z/n$ is generated by $g=A_s^{-1}e_2$, $a=a_2+a_{\rm odd}$ with $a_2=(\tfrac12,0)$ and
$a_{\rm odd}=(1,-1)/(s-2)\in L'$, and in the basis $u'_{x,\epsilon,\delta}:=u_{xg+\epsilon a_2,\,-\delta a_2}$ ($N_s=4n$ labels):
\[
T=\mathrm{diag}_x\,\e(-x^2/n)\otimes T_2,\qquad S_{\rm lab}=\bigl[\e(-2xy/n)/\sqrt n\bigr]_{x,y}\otimes S_2,
\]
with $S_2:(\epsilon,\delta)\mapsto(\delta,\epsilon)$ with phase $(-1)^{\epsilon\delta}$, and $T_2$ as above; $S_2,T_2$ do not depend on $s$.
(i) With $\rho(S):=S_{\rm lab}^{-1}$, $\rho=\bigotimes_{p^e\Vert n}\rho_{q_p}\otimes\rho_2$, where $\rho_{q_p}$ is the Weil representation of $(\Z/p^e,q_p)$,
$q_p(x)=-v_px^2/p^e$, $v_p=(n/p^e)^{-1}\bmod p^e$, and $\rho_2=(\e(3/4)S_2^{-1},T_2)$ is the $4$-dimensional $s=7$ factor. The Weil scalar is $\e(1/4)$ for every odd $s$, since
$n\equiv3\bmod4$. $\rho$ factors through $\mathrm{SL}_2(\Z/4n)=\mathrm{SL}_2(\Z/4)\times\prod\mathrm{SL}_2(\Z/p^e)$, so $N_s=4s(s-2)$ is the number of labels and the order of $T$.
(ii) Negation acts as $\bigotimes_pZ_p\otimes I_4$; $\gamma$ acts as $x\mapsto(s-1)x$ with phase $(-1)^{\epsilon\delta}$, i.e.\ as
$\bigotimes_{p\mid s}Z_p\otimes\bigotimes_{p\mid s-2}I\otimes Z_2$. In fact $\gamma\equiv(s-1)\cdot I$ on all of $L'/\Z^2$, because
$(\gamma-(s-1)I)A_s^{-1}=\begin{pmatrix}0&-1\\1&0\end{pmatrix}$; and $s-1\equiv-1\bmod p^e$ for $p\mid s$, $\equiv+1$ for $p\mid s-2$ (the reduction of
$\varepsilon_s=(s-1)+\sqrt{n_s}$ at the ramified primes).
(iii) $P=S_{\rm lab}^2=Z_n\otimes I_4$, so $\dim W_P=4\cdot\tfrac{n+1}2=2(s-1)^2$ (this is $72$ at $s=7$).
(iv) $\mathrm{Fix}(N)\cap\mathrm{Fix}(\gamma)=\{Z_{(s)}Z_{(s-2)}=+1,\ Z_{(s)}Z_2=+1\}$ has dimension $3E(s)E(s-2)+O(s)O(s-2)=s(s-1)$, with $E(m)=\tfrac{m+1}2$,
$O(m)=\tfrac{m-1}2$ and $Z_{(m)}=\bigotimes_{p\mid m}Z_p$. Hence $\dim V\le s(s-1)$. There are exactly $2s-3$ forced-zero labels, and the odd-characteristic
sector $B_1\otimes(\Z/s)^-\otimes(\Z/(s-2))^-$ has dimension $(s-1)(s-3)/4$.
\end{theorem}
\catA{Symbolic proof for all odd $s$ (the $T$- and $S$-entries from $B(g,g)=-2/n$, $B(g,a_2)=0$, $w_c=g/2$, and the CRT split of $-x^2/n$); exact verification
entry by entry at $s=5,7,9,11,13,15,21,25$ (including the prime powers $9$ and $25$). The congruence statement in (i) is conditional on the cited
theorem \cite{Weil}; its convention was checked against the exact $s=7$ relations. Note $\dim W_P=2(s-1)^2$, while $3\prod(p^e+1)/2$ is the bound for $d_{\rm Eis}$.}

\begin{remark}[The two generators at $s=7$]\label{rem:gen}
Theorem~\ref{thm:tensor7} uses $g=(7,3)/35$, Theorem~\ref{thm:tensorodd} uses $A_7^{-1}e_2=(7,-2)/35=11\cdot(7,3)/35$ in $L'/\Z^2$. By CRT the
local factors of Theorem~\ref{thm:tensor7}(i) combine to the $\Z/35$-phases $\e(13xy/35)/\sqrt{35}$ for $S_{\rm lab}$ ($\tfrac45+\tfrac47=\tfrac{48}{35}$) and
$\e(24x^2/35)$ for $T$ ($\tfrac25+\tfrac27=\tfrac{24}{35}$). Under $x\mapsto11x$ these become $\e(-2xy/35)$ and $\e(-x^2/35)$, since $121\cdot13\equiv-2$ and
$121\cdot24\equiv-1\pmod{35}$; these are the phases of Theorem~\ref{thm:tensorodd} at $n=35$. The two descriptions are the same representation in two bases.
\end{remark}

\subsection{Sectors: parity and the orbit lemma}\label{sec:sectors}
\begin{proposition}[Odd-parity sectors are cusp]\label{prop:oddcusp5}
If $\rho(T)$ has no eigenvalue $1$ on a $\rho$-invariant subspace $\sigma$, the projection of the lift onto $\sigma$ is cuspidal. Consequently $d_{\rm cusp}\ge(s-1)(s-3)/4$ wherever the $B_1$ sector lies in $V$, and
$d_{\rm Eis}\le3\prod(p^e+1)/2$ whenever no odd parity is carried only by primes with $e\ge2$ (every tested odd $s\le99$; the first undecided cases are $s=225,1225$).
\end{proposition}
\catA{Proved: no constant term at any cusp. \emph{Category~B:} equality $d_{\rm Eis}=3\prod(p^e+1)/2$, i.e.\ that the all-even sector carries no cusp content.}
\begin{proposition}[Parity fact]\label{prop:parity5}
For every odd $s$, every $\Delta$ label has even second coordinate $x_2$. Hence $Z_{\rm all}$ acts as negation on $\Delta$, fixes every $\langle\gamma,-1\rangle$-orbit, and every sector with $\varepsilon_{\rm all}=-1$ vanishes.
\end{proposition}
\catA{Symbolic: $D\cdot A^{-1}=-2\,\mathrm{adj}A$ gives even second coordinates on $L'$; $a=(s^2,-2s)$; $D$ is even.}
\begin{proposition}[Orbit lemma]\label{prop:orbit5}
For every odd $s$, $c_\varepsilon(O_0,q_{\min})=(1+\varepsilon_{\rm all})N_{O_0}(q_{\min})$, where $O_0=\{\pm\nu_0,\pm\gamma\nu_0\}$ and $\nu_0=(s,-2)$. This is nonzero on every sector with $\varepsilon_{\rm all}=+1$. Extra reduced classes do occur, at $s=27,245,363,847,2025,\dots$ ($16$ cases with $s\le30001$), but they never land in $Z_SO_0$.
\end{proposition}
\catA{Symbolic (Nagell bound, $\gamma^2\nu_0\equiv\nu_0$, stabiliser $\{1,Z_{\rm all}\}$, exclusion of extra classes). Also checked exactly for odd $s\le401$ and at all $16$ exceptions. At $s=27$ the sector run certifies $\dim V=702$ (exact), in agreement with Theorem~\ref{thm:dims}.
Combined with Proposition~\ref{prop:parity5}, the matched sectors with $\varepsilon_{\rm all}=+1$ have a nonzero lift. Injectivity therefore follows wherever the sector is irreducible (squarefree $s(s-2)$). \emph{Open:} $\dim V=s(s-1)$ when $e\ge2$.}
